% Detailed modem view. Included from the main whitepaper; not a standalone paper.
\section{Agent orchestration and modem software architecture}
\label{sec:modem-software}

This section connects the proposed ports-and-adapters architecture to the
Rust modem implementation as of 11 November 2026. The radio-free modem
preview serves traditional clients over paired TCP sockets and exchanges
audio with a bench peer. The proposed station runtime adds policy, device
ownership, and physical transmit control. The evidence table records
modem interoperability, client transfer, and station qualification separately.

\subsection{Standalone modem product and shared runtime}
\label{sec:modem-shared-runtime}
\emph{Proposed integration.} The native modem should be delivered as a shared
Rust runtime usable by traditional amateur-radio applications without an AI
host. The foreground executable \path{linkwise-modemd} will
assemble the runtime and its adapters. Optional operating-system supervision
and graphical configuration do not define a different protocol engine.
Linkwise orchestration remains an additional client of the station execution
contract, not a prerequisite for an ordinary client connection.

The common core owns session state, queues, deadlines, handshake and ARQ
sequencing, byte accounting, and outcomes. Platform adapters provide audio
capture and playback, device discovery, PTT, permissions, and process startup.
The implementation supports virtual time for deterministic tests and a
monotonic clock for the foreground runner. The device execution tests use
simulated audio and PTT.

The radio-free preview supports uncompressed BW500 and fixed level-4
BW2300/BW2750. BW500 remains the first target for station integration.
A release manifest must enumerate the qualified roles, DATA forms,
host commands, reply/query forms, rate behavior, and disconnect behavior.

\subsection{Traditional-client interface and session admission}
\label{sec:modem-client-admission}
\emph{Proposed integration.} A VARA-compatible host adapter exposes a
CR-delimited command/notification TCP channel and a separate raw binary DATA
channel. Ports default to loopback 8300/8301 where available and remain
configurable. These are host-protocol compatibility goals; neither accepting
a command nor opening a socket establishes on-air interoperability. Winlink
message framing, mailbox exchange, and message acceptance remain the client's
responsibility. No message body or continuous modem stream traverses MCP.

The preview uses \path{host::tcp} to pair one host command socket and one
DATA socket. \path{api::CommandDecoder} parses the command stream.
A separate IPv4 loopback socket exchanges audio with a single cable peer.
Each record has a four-byte big-endian sample count followed by little-endian
32-bit float samples at 48 kHz mono. The audio bridge can attach named
virtual capture and playback devices.
A compatibility matrix must record initialization commands, replies,
asynchronous notifications, buffer semantics, and unsupported behavior for
each qualified client/version. The current parser supports only a subset;
unsupported commands continue to fail explicitly. In particular, an
\texttt{OK} reply acknowledges command acceptance, not a completed RF link.

One station execution service arbitrates both admission paths:
\begin{description}
\item[Standalone.] A locally configured client principal requests
\texttt{CONNECT}. The service validates identity, station configuration,
route readiness, lease availability, and installed authority, then creates
one bounded operation and native session. Listening also requires an installed
receive/respond policy; a decoded inbound call does not itself authorize PTT.
\item[Managed.] The coordinator prepares and authorizes an attempt and
reserves an expected client association. The accepted start arms that
association; the matching client \texttt{CONNECT} consumes it once and starts
the backend. It must not create a second operation. The reservation binds
client ownership, local/remote callsigns, configuration revision, lease, and
expiry. Mismatches, competing requests, and attachment timeout fail explicitly.
\end{description}

The legacy two-socket protocol supplies no authenticated pairing token.
Initially permit one exclusive control/data pair per configured endpoint,
with a bounded attachment timeout and an explicit pairing state machine that
handles either arrival order. RF admission requires a complete pair.
A replacement socket cannot adopt an existing
session; when either channel closes, both are detached, admission is blocked
until cleanup completes, and the active session is cancelled with PTT release
requested. Peer address alone is not client identity. Loopback is a deployment
restriction rather than authentication; shared or remote deployments require
an authenticated enclosing transport or a separately qualified ownership
mechanism. Model disconnect is different from modem-client channel loss.

\subsection{Native session and device execution contracts}
\label{sec:modem-execution-contracts}
\emph{Proposed integration.} One serialized native session actor owns the
current peer, protocol phase, transmit/receive epoch, outstanding DATA, and
monotonic deadlines. Its lifecycle distinguishes idle, admitting,
establishing, connected, closing, completed, aborted, and faulted states;
transmit, answer-wait, query, and turn phases are subordinate protocol states.
Only validated events in the current phase may advance it. Every asynchronous
audio or device result carries a session identity and a unique transmission
token; results from a previous session or transmission are ignored and
recorded. Stop and channel-loss cleanup are idempotent. An accepted abort
invalidates outstanding work before any late completion can advance state.

Application lifecycle requests enter through \path{ModemSessionPort}; host
commands enter through \path{host::ModemPort} and the same admission service.
The native backend adapter binds both to the one actor. It must not run a
second protocol machine in the TCP handler or in the station coordinator.
External VARA and simulation remain alternative backends behind the lifecycle
port; only the native backend assembles the Rust DSP/session components.

Audio contracts identify a continuous capture epoch, absolute sample ranges,
qualified route and format, and terminal discontinuity. Playback accepts
owned bounded buffers and reports accepted, started, presented, cancelled,
or failed with the transmission token. Queue acceptance or callback buffer
consumption is not presentation completion. Cancellation must establish
whether queued samples can still reach the device; uncertainty faults the
session and prevents a new transmission. The PTT contract separates permit,
assertion, release request, and observed release. Actual presentation plus
qualified release/turnaround evidence is required before the actor emits a
protocol transmit-complete event and arms its receive deadline. Independent
watchdog release does not wait for the actor, event storage, or an AI host.

All mailbox, audio, event, and byte queues have configured finite limits.
Backpressure suspends acceptance at the host boundary without blocking DSP
callbacks or the actor's stop/watchdog processing. Protocol deadlines cannot
be extended by a slow client or a full notification queue. A stalled consumer
has a bounded timeout and a terminal reason; events needed to account for
bytes cannot silently disappear. Device loss, capture gaps, lease expiry,
playback timeout, and unconfirmed PTT release terminate the normal path.

Byte accounting distinguishes cumulative accepted bytes, pending bytes,
unique outstanding bytes, peer-acknowledged bytes, unresolved bytes, and
explicitly discarded bytes. For an operation, accepted bytes equal pending
plus outstanding plus acknowledged plus unresolved plus discarded bytes;
a byte occupies only one of those categories at a snapshot. Transmit attempts
and retries are separate counters and never inflate acknowledged totals.
Receive accounting distinguishes bytes decoded, reassembled/deduplicated,
queued for the host, and written to its socket. Socket write completion is
not proof the client consumed or committed the bytes. Host \texttt{BUFFER}
notifications need a client-qualified mapping to these counters rather than
the current ARQ queue size. Peer acknowledgment is not Winlink message or
CMS acceptance.

On ambiguous failure the operation result records unresolved byte counts and
ranges; payload remains in bounded session-owned memory rather than audit
events. Explicit abort/clear can discard bytes with a recorded count. No reconnect
or process restart replays them automatically. In-memory retention does not
promise crash durability: a crash produces an unknown outcome unless a
separately implemented bounded journal supplies evidence. Native on-air
duplicate handling must be qualified independently of local callback tokens.

\subsection{Shared-core implementation order and exit criteria}
\label{sec:modem-core-plan}
The first two stages have an implementation at the validated-event boundary
with a clock-driven runner and simulated devices. The BW500 audio loopback
and paired TCP bench exercise stages three and four with simulated peers.
Rolling pilot acquisition for levels 1--3 has separate off-air receive
evidence (Table~\ref{tab:modem-evidence-gates}). The remaining work includes
production DATA sequencing and recovery, native audio/PTT integration,
and broader client and reference-peer qualification. The sequence below defines the
scope and exit criteria for each stage.
\begin{enumerate}
\item \textbf{Session contracts and owner.} Define typed commands/events,
state transitions, operation/session/transmission identities, finite budgets,
byte-accounting rules, and a virtual clock. Demonstrate rejection of competing
owners, illegal transitions, stale completion, and repeated cleanup.
\item \textbf{Simulated audio and PTT execution.} Bind capture epochs,
playback tokens, presentation, release, and turnaround to the actor. Inject
late/missing completion, capture gaps, cancellation races, device loss, and
watchdog expiry. Demonstrate independent release and no false TX completion.
This yields the first foreground shared-core milestone without native devices.
\item \textbf{First complete link profile.} Integrate establishment,
packetization, DATA fields, validated replies/queries, ARQ, duplicate handling,
turn/rate ownership, keepalive, and disconnect. Exchange arbitrary binary
streams between two runtimes under deterministic loss, duplication, and delay;
preserve byte-accounting invariants and identify unresolved outcomes. Resolve
wire uncertainties with independent evidence rather than invented sequencing.
\item \textbf{Traditional-client service.} Add loopback listeners, exclusive
pairing, admission, notifications, timeouts, and bounded backpressure. A host
harness must exercise initialization, both socket arrival orders, competing
clients, binary preservation, stalled consumers, channel loss, and reconnect.
Only advertise the tested host-protocol subset.
\item \textbf{Station integration.} Bind standalone and managed admission to
policy, leases, grants, reservations, and operation/event records. Prove a
matching managed \texttt{CONNECT} starts exactly one attempt and a mismatch
cannot transmit. Test client loss separately from model loss.
\item \textbf{Reference and client qualification.} Demonstrate bidirectional
transfers with an unmodified release-specific reference modem and a named
Winlink client/version. Preserve commands, both audio directions, timing,
payload comparisons, recovery evidence, and distinct modem/client outcomes.
Native hardware measurements still require the corresponding device adapter.
\end{enumerate}

Only after these shared contracts stabilize should platform work supply
macOS audio/permissions and launch-agent integration, Linux audio/device
permissions and user-service integration, and Windows audio/device access
and background-process packaging. Foreground operation comes first on each
platform. A system service is optional and requires qualification in its
actual execution context; it is not evidence of audio readiness. Each native
adapter must satisfy the same fault/completion contract and undergo device
and inhibited/conducted PTT qualification before operational RF use.

\subsection{ChatGPT Agent and Codex as the orchestrating agents}
ChatGPT Agent or Codex is the decision-making orchestrator: it observes station
readiness, compares eligible RMS channels and evidence, chooses the next
attempt, invokes MCP tools, and adapts its plan to recorded outcomes. These
are alternative AI-host clients of the same typed application contract;
client transport and reconnect behavior must be qualified for each deployed
host. Neither client is a modem timing source, a radio driver, or an authority
issuer. The agent's context and conversation history are not station state.

The local link coordinator is the \emph{station execution service}. It performs
bounded use cases under deterministic policy rather than operating a second,
competing strategic planner. The agent may choose a candidate; the domain
ranker supplies reproducible scores and exclusions; local policy decides
eligibility and authorization; the execution service serializes preparation
and supervises the accepted session. Handshake, ACK/NAK, retransmission,
audio-buffer, and PTT deadlines stay local even while the agent is thinking,
disconnected, or recovering context.

The preparation contract remains the one in the interface section:
\path{link_prepare_candidate} returns a short-lived handle bound to candidate,
lease, verified configuration revision, route profile, and expiry;
\path{link_request_session_authorization} asks deterministic policy for a
separate bounded grant; \path{link_start_session} consumes both after a fresh
validation. Authorization is automatic within installed station authority,
not a runtime human approval step. Every mutation carries principal,
request/correlation IDs, policy revision, deadline, and an idempotency key.
Repeating an accepted start request must return the same operation rather than
create another RF attempt.

An accepted start returns a durable operation ID and session identity promptly.
The agent follows bounded session snapshots and ordered cursor-based events;
it does not hold one MCP request open for an entire contact. Durable operation
records, terminal reasons, and explicit unknown/stale fields let a reconnecting
agent reconstruct what happened. MCP disconnect and application restart are
different events: local policy determines whether an already accepted session
finishes or aborts after client loss, while process restart invalidates
in-memory grants/leases and requires reconciliation. Neither event extends a
transmit permit or authorizes automatic replay of unresolved bytes.

\subsection{Ports, ownership, and backend interchangeability}
The inbound MCP adapter validates identity and request shape and translates
results. Application services own preparation, deadlines, operation records,
and lifecycle orchestration. Domain services own policy, leases, session
invariants, and authorization. Concrete adapters own device/protocol details.
The proposed modem backend contract is \path{ModemSessionPort}; the existing
\path{host::ModemPort} is an inbound host-command/data boundary and must not be
mistaken for this outbound driver interface.

\begin{longtable}{@{}p{0.25\textwidth}p{0.69\textwidth}@{}}
\caption{Ownership at the agent, station, and modem boundaries.}\label{tab:modem-ownership}\\
\toprule
\textbf{Owner} & \textbf{Responsibility and boundary}\\
\midrule
\endhead
AI host & Select objectives and subsequent attempts from observed evidence;
request typed actions without minting grants, changing installed limits, or
addressing hardware sockets.\\
Station execution service & Own attempt/session identity, prepared handles,
lease adoption, client/modem association, cancellation, reconciliation, and
terminal operation results.\\
Modem backend adapter & Report capabilities and readiness; accept internal
bounded lifecycle requests; expose authoritative protocol events and a
backpressured data endpoint. Implementations may be native Rust, qualified
external VARA, or simulation.\\
Native session owner & Own one protocol state machine, receive/transmit epoch,
current peer, timers, outstanding bytes, and the ordering of DSP and playback
work. The simulated runtime implements session ownership and deadlines.\\
Audio adapter & Own configured input/output routes, formats, buffers, sample
continuity, and presentation-completion evidence. CPAL provides input capture and an output
primitive with a presentation estimate. Hardware completion remains open.\\
PTT arbiter & Own the only software key-up path, one-use permits, route checks,
key-up and session budgets, watchdog release, and requested/observed state.\\
Winlink client adapter & Own its supported queue, message-service session,
transfer semantics, and confirmation evidence. It does not delegate message
framing or body bytes to the language model.\\
\bottomrule
\end{longtable}

\landscapeplantuml{build/diagrams/Linkwise_Modem_Runtime}{Proposed agent-to-modem ownership and backend choices, over the native components currently implemented. Dotted integration edges remain incomplete; this is not an as-built operational runtime.}{fig:modem-runtime}

The backend port should provide versioned capabilities, an observed state
snapshot, internal begin/stop operations, a session/event cursor, and a bounded
data-stream attachment. The native adapter and external VARA adapter must
produce the same domain event classes while preserving native event codes,
backend versions, units, timestamps, unsupported fields, and qualification
provenance. A simulator implements that same port without exposing hidden RF
truth through MCP. Backend replacement cannot silently increase station
privileges or normalize an unverified signal into a connection.

\subsection{MCP observations and the modem data plane}
The modem's MCP surface is observational:
\path{modem_get_capabilities}, \path{modem_get_state},
\path{modem_get_session}, and \path{modem_get_events}, plus the cached
\path{linkwise://modem/state} resource. Session start/abort is routed through
the coordinator, and explicit CAT/PTT tools remain available through their
policy-mediated contracts. Audio-route selection uses qualified profile IDs,
never arbitrary device names, filesystem paths, gain scripts, or routing
expressions supplied by the agent.

Capability results must distinguish implementation from qualification: for
example, level-4 single-frame synthesis exists, but stock-peer acceptance and
multi-frame operation are unqualified. An advertised capability needs a
backend/version, mode, direction, supported parameters, availability, and a
qualification scope. State should distinguish device present, stream open,
audio observed, synchronized candidate, CRC-valid frame, handshake complete,
and active data exchange. Observations include sample ranges where useful and
source/freshness metadata; missing BER, SNR, throughput, or retry measurements
remain unsupported, not zero.

Protocol bytes travel between the Winlink client and modem on a bounded
stream or spool. MCP carries aggregate queue counts, message references where
supported, progress, and outcome summaries; raw audio, continuous IQ, and
message bodies are excluded by default. The agent-facing operation ID, station
attempt ID, modem-session ID, and Winlink-session ID remain distinct and are
correlated in the event store. UTC timestamps support audit correlation;
monotonic instants enforce deadlines; absolute sample counters locate signal
evidence. One must not substitute for another.

\subsection{Native implementation map and current limits}
The Rust package contains the foreground bench service and reusable components.
\path{audio::device} captures input; \path{modem::receiver} bounds and dispatches
receive candidates; \path{modem::waveform} supplies energy detection and tone
correlation; \path{modem::bw500} implements index-modulated levels 1--3;
\path{modem::bw500_level4} handles level-4 channel coding;
\path{modem::bw500_waveform} adds measured-reference two-band synthesis and
reception; \path{modem::wide_level4} supplies level-4 setup and level-2 through
level-4 DATA waveforms for BW2300 and BW2750, using the carrier allocations
in \path{modem::wide_allocations}. \path{modem::coding} implements whitening, CRC, interleavers, and
FEC. \path{modem::control}, \path{modem::setup}, \path{modem::handshake}, and
\path{modem::audio_handshake} cover keyed control, caller setup, and
establishment. \path{modem::data} and \path{modem::arq} are frame/container
and sender components. \path{api::CommandDecoder} and \path{host} implement
host stream boundaries.

The clock-driven \path{modem::runtime} combines establishment and ARQ at
the validated-event boundary. The virtual host in \path{linkwise-modemd}
owns its bench handshake, transmit turn, pending DATA, and deadlines.
The host keeps command and binary DATA streams on paired TCP sockets. A
separate virtual audio connection carries a big-endian sample count followed
by little-endian 32-bit floating-point samples at 48 kHz. The host queues
up to 4,096 bytes and segments BW500 writes into measured uncompressed
deliveries of up to 203 bytes. A wide level-4 delivery can carry two 89-byte
full frames and a closing frame of up to 86 bytes. Wide level-3 and level-2
full frames carry 47 and 22 bytes, respectively.
The optional CPAL bridge provides native capture and playback for isolated
audio routes. Station integration still needs qualified hardware playback
completion, PTT release, and the station admission and ownership contracts.

ARM64 preview installers and station guides are available for macOS,
Windows, Ubuntu, and Fedora. Windows, Ubuntu, and Fedora passed install,
launch, simulation, and removal checks. The extracted macOS package
executable ran successfully, and its bundled guide was checked. Native macOS
capture awaits an input endpoint.

\section{DSP implementation pipeline}
\label{sec:modem-dsp}

The standalone modem implementation reference expands the implemented
waveforms, trellis metrics, and bench recovery traces.
Its source and build instructions are in \path{docs/modem-implementation/}.

% Shared by the architecture white paper and the modem implementation reference.
\providecommand{\modemcite}[1]{}

\subsection{Audio ingress, continuity, and candidate detection}
The optional CPAL input backend requires native 48,000 samples/s and a selected
zero-based channel; it accepts float32, signed 16-bit, or unsigned 16-bit
formats and converts to normalized mono float samples. It selects one channel
rather than averaging potentially phase-opposed channels. No implicit
resampling is implemented. The callback accumulates 1,024-sample blocks and
uses a nonblocking bounded queue of 64 blocks. This represents approximately
1.365 s of mono input; it is a queue capacity, not a permitted DSP or keying
latency. Queue overflow, a malformed interleaved input block, or device failure
latches a terminal capture fault. The receiver must restart on a new continuous
stream instead of concealing dropped samples.

DSP runs outside the capture callback. The streaming receive worker validates
finite samples before advancing state, bounds processing steps to 480 samples,
and retains at most 30 s of mono history (1,440,000 float samples, about
5.76 MB of sample payload). Demodulation allocates additional bounded working
arrays, so the retained-history limit is not a whole-process memory limit.
Absolute half-open sample ranges survive chunk boundaries and ring-buffer
wrap. If a candidate began before retained history, it is marked truncated and
emits no partial frame. WAV reports count truncated energy candidates and
retain independently acquired pilot frames.

The energy detector accumulates 10 ms windows. A window is active when
\[
 \frac{1}{W}\sum_{n=0}^{W-1}x[n]^2 \geq 10^{-4},\qquad W=480.
\]
A 10 ms window holds 480 samples at 48 kHz. Square each normalized
sample so positive and negative excursions both count, then average.
A mean square of $10^{-4}$ corresponds to an RMS amplitude of 0.01,
so this gate opens at one percent of full scale.

A quiet window closes a candidate; end-of-stream flushes a partial window and
any open candidate. The 1\% full-scale RMS gate serves level-4 acquisition
with an external measurement profile and reports energy candidates for
diagnostics. The CLI scans level-1--3 pilots independently of this gate.
Energy candidates can include impulsive noise and unrelated signals.

\subsection{BW500 levels 1--3: tone-index geometry}
\modemcite{proakis}These levels encode data by selecting a frequency bin for each symbol. They
are not the level-4 differential-BPSK waveform played more slowly. The
implementation performs direct sine/cosine dot products at selected DFT bins;
it does not compute a full FFT for each symbol.

\begin{table}[htbp]
\centering\small
\begin{tabular}{@{}lrrrrrrr@{}}
\toprule
Level & $H$ & Bins & Columns & Refs. & $N$ & Coded & Payload\\
\midrule
1 & 2048 & 54--74 & 124 & 24 & 96 & 300 & 9 B\\
2 & 1024 & 27--37 & 228 & 24 & 200 & 612 & 22 B\\
3 & 1024 & 27--37 & 226 & 24 & 297 & 606 & 34 B\\
\bottomrule
\end{tabular}
\caption{Implemented index-mode geometry at 48 kHz. Columns exclude the initial lead; $N$ includes mode padding before whitening. Payload excludes the control byte and two-byte CRC.}
\label{tab:dsp-index-geometry}
\end{table}

The symbol rates are $48000/H$: 23.4375 symbols/s at level 1 and 46.875 at
levels 2/3. The synthetic transmitter adds two lead symbols at level 1 and four
at levels 2/3, all totaling 4,096 samples. Those bench leads use fixed in-band
tones; they do not implement the stock session's randomized lead stream.
There are three coded bits per data column. Subtracting the 24 reference
columns gives $100\times3=300$, $204\times3=612$, and $202\times3=606$
channel bits respectively.

Mode construction uses seeds 12,600,710; 7,910,701; and 14,955,604. One generator
stream produces the channel permutation, then the turbo permutation, followed
by bin/pilot draws. Level 1 replaces its generated base bins with the explicit
measured 124-entry array in the source. Levels 2/3 currently use generated
base allocations. All three therefore need qualification appropriate to their
actual construction, not a blanket claim that every allocation is independently
measured.

For each data column, the eight labels $(0,1,3,2,7,6,4,5)$ select offsets from
the column's base bin. Level 1 applies stride 2 modulo a span of 21; levels
2/3 apply stride 1 modulo 11. Fixed reference columns and their known bin
values allow synchronization without comparing decoded payload to an expected
message. The level-3 information block contains 37 complete bytes plus one
padding bit: 34 payload bytes, one control byte, two CRC bytes, and a zero
padding bit before whitening.

\landscapeplantuml{build/diagrams/Linkwise_DSP_Index}{Implemented BW500 level 1--3 receive pipeline. Pilot acquisition locates index frames in continuous audio. Soft demodulation and CRC validation follow timing and frequency selection.}{fig:dsp-index}

\subsection{Index-mode acquisition and soft demodulation}
\modemcite{proakis,meyr}The WAV and device receive commands use \path{modem::bw500::scan} to
search every 128 samples in a bounded rolling spectrum history. At each
epoch, the scanner scores the 24 reference columns against energy in
the occupied band and guard bins at offsets of 0 and $\pm$12 Hz. It
considers hypotheses with a mean pilot score of at least 0.65 and
refines timing in 16-sample steps. This lets a frame span brief drops in
energy or continuous receiver noise. The library also retains the
candidate-based decoder, which starts its search near $b+LH$ for burst
start $b$ and lead length $L$. The score threshold has no measured
false-alarm rate.

For a symbol of length $H$ and bin $k$, the unnormalized tone energy is
\[
 E_k=\left|\sum_{n=0}^{H-1}x[n]\exp(-j2\pi kn/H)\right|^2.
\]
The exponential supplies a reference that completes $k$ cycles over the
$H$ samples. Multiplying and summing aligns a matching tone so its
contributions add together. Taking the squared magnitude measures the
strength of that match independently of its starting phase. For example,
at 48 kHz with $H=1024$, bin 32 tests for a 1,500 Hz tone.

For timing refinement and soft decisions, the implementation divides by
$(H/2)\sum x[n]^2$, with an epsilon floor. For data, it measures the
eight mapped bin hypotheses. For each of the three bit positions, the soft
metric is the largest normalized power among labels with bit 1 minus the
largest among labels with bit 0. This max-based difference is a decoder input;
it is not a calibrated log likelihood, SNR, or BER estimate.
A positive value means the strongest tone compatible with bit 1 beat
the strongest tone compatible with bit 0. Values near zero give the
decoder little reason to prefer either bit.

Metrics are scattered through the inverse channel permutation into FEC order.
Only a CRC-valid result is returned. The rolling scanner suppresses overlap
with an emitted frame and continues to search for adjacent frames. Level 4
handles one frame per energy candidate. Continuous carrier and sample-clock
tracking remain open.

\subsection{Index-mode transmit synthesis}
The level-1/2/3 transmitter accepts a nonempty list of frames with exactly
9, 22, or 34 payload bytes per frame. It appends the control byte and byte CRC,
serializes MSB-first, pads to the mode's $N$, whitens, and turbo-encodes. For
each data column it gathers three channel-ordered coded bits into an integer
0--7 and selects the corresponding mapped tone. Reference columns emit their
fixed expected bins without consuming coded bits. A burst has one lead at its
start, followed by consecutive frame-column blocks; the lead is not repeated
between frames.

Each selected bin $k$ is rendered as $0.5\sin(2\pi kn/H)$ for $n=0,\ldots,H-1$.
The phase argument restarts for each symbol; integer-bin tones complete an
integer number of cycles over $H$ samples. The implementation concatenates
those symbol buffers directly. It does not add DATA-symbol pulse shaping,
resampling, a stock randomized acquisition stream, or a qualified transmitter
spectral mask. The fixed half-scale amplitude is a bench sample convention.
Reference-peer acceptance, occupied bandwidth, audio-route gain, and actual
RF power must be measured separately before treating this synthesizer as an
operational transmitter.

\subsection{Shared CRC, whitening, and interleaver generation}
\modemcite{koopman}Frame integrity uses CRC-16/GENIBUS: polynomial \texttt{0x1021}, initial register
\texttt{0xffff}, final complement, no reflection, and a big-endian two-byte
trailer. The check vector is \texttt{123456789} $\mapsto$ \texttt{0xd64e}.
The CRC covers the actual frame body, including its control field; mode padding
bits are outside the byte CRC. CRC checks establish frame integrity under this
code, not authenticated sender identity or recipient delivery.

Bits are serialized most-significant first, padded to the mode's information
size, and XORed with a freshly seeded whitening sequence. The 24-bit generator
obeys
\[
 s_{i+1}=(\texttt{0x43fd43fd}\,s_i+\texttt{0xc39ec3})\bmod 2^{24}.
\]
Each step multiplies the current state $s_i$ by the first constant,
adds the second constant, and keeps the lowest 24 bits. The remainder
operation is what makes the state wrap. The same starting state and
steps produce the same bit sequence at both ends of the link.

Whitening starts at $s_0=29318$ and takes the next state's high bit. A
generator bit of 1 flips the corresponding frame bit. A generator bit
of 0 leaves it alone. Applying the same mask again flips those bits
back. Permutations use
rejection sampling: draw an index, append only its first occurrence, and
continue until every index is present. The draw multiplies the normalized
state by the requested size with single-precision rounding, matching the
implementation's wire-oriented generator semantics. Changing numeric precision
or substituting a library shuffle can change the permutation.

The channel permutation and turbo permutation are different objects. If
$\pi_1$ maps transmitted position $j$ to a coded-bit position, TX gathers
$q_j=c_{\pi_1(j)}$ and RX scatters the corresponding soft metric back into
$c_{\pi_1(j)}$. The inner permutation $\pi_2$ changes the input order of the
second constituent encoder and must be inverted consistently when exchanging
extrinsic information. Interleavers belong to the mode/profile, not to the
agent's request parameters.
In the worked example, the channel permutation puts coded bit 33 in the
first transmitted position. The receiver puts the first soft decision
back into slot 33 before decoding. The turbo permutation acts earlier,
when the second encoder reads the information bits in a different order.

\subsection{Turbo encoder and iterative receive decoder}
\modemcite{berrou}The codec has two eight-state recursive systematic convolutional constituents.
Writing the three state bits as $(r_0,r_1,r_2)$ and the next information bit as
$u$, the implementation uses
\[
 a=u\oplus r_1\oplus r_2,\quad
 p=a\oplus r_0\oplus r_2,\quad
 (r_0',r_1',r_2')=(a,r_0,r_1).
\]
The symbol $\oplus$ means exclusive OR. First compute the new feedback
bit $a$ from the input and two stored bits. Then compute the parity bit
$p$ from $a$ and two stored bits. Finally, shift the memory so $a$
becomes the newest bit. Starting with an empty register and input 1,
the encoder outputs parity 1 and moves to state $(1,0,0)$.

Both encoders start at zero. The first consumes natural-order bits and the
second consumes turbo-permuted bits. Levels 1/2 transmit systematic and both
parity streams (rate 1/3); levels 3/4 transmit systematic bits and alternate
which constituent parity is retained (rate 1/2). Each constituent is terminated
with three steps using $u=r_1\oplus r_2$. The wire tail is twelve bits,
interleaving the two systematic/parity tail pairs at each of three steps.
Thus the implemented coded sizes are $3N+12$ or $2N+12$.
The factor of three or two counts the transmitted bits per information
bit before termination. For the worked level-4 frame, 368 information
bits give 736 systematic and parity bits, then the tail brings the
total to 748.

\landscapeplantuml{build/diagrams/Linkwise_DSP_FEC}{Implemented shared whitening, constituent coding, tail and permutation structure. Receive decoding uses a hard-decision CRC fast path followed by at most eight max-log-MAP iterations.}{fig:dsp-fec}

\modemcite{bcjr,robertson}The receiver first makes hard decisions on systematic metrics and checks the
frame CRC. If that fails, it constructs the two terminated trellises, treating
punctured parity as a zero-information metric. A forward/backward
\emph{max-log-MAP} pass computes extrinsic information for each constituent.
The implementation uses maxima rather than log-sum-exp correction; it must
not be described as an exact log-MAP decoder. Positive input metrics favor
bit 1. The two constituents exchange interleaved/de-interleaved extrinsic
information for at most eight iterations, checking de-whitened frame integrity
after each iteration and stopping on a valid CRC.

This implementation does not maintain cross-retransmission soft combining.
Its metrics have heuristic scales and its tests do not establish coding gain,
waterfall curves, or residual undetected-error probability under HF fading.
Those require independently generated vectors and controlled channel sweeps,
including offset, burst noise, clipping, and long erasures.

\subsection{Level-4 frame geometry and external measurement profile}
The level-4 byte container is 46 bytes: 43 payload bytes, one control byte,
and two CRC bytes. Its $N=368$ information bits produce 748 coded bits. The
channel and turbo permutations start from seed 6,437,280. Pilot generation
continues that stream after the two permutations, choosing 18 distinct columns
per band with $5+\lfloor R\,386\rfloor$.

The implemented waveform stores 393 columns per band: one reference followed
by 392 body columns. Removing 18 pilot positions in each band yields
$2(392-18)=748$ differential data cells. This is the retained single-frame
implementation geometry, not a claim that every stock burst has exactly that
many emission slots. Its reference centers start after nine 512-sample slots.
Band centers are 1,350 and 1,650 Hz, the symbol hop is $H=512$ samples
(93.75 symbols/s), and the upper band is staggered by 256 samples. Unlike the
index ladder, these carriers are not selected integer bins of a symbol FFT.

The phase reference consists of 786 measured points on a 128-phase lattice,
393 per band. Each point is $G_{b,t}=\exp(j2\pi k_{b,t}/128)$. A shared measured
513-sample real pulse covers offsets $-256$ through $+256$. The profile loader
requires all reference positions, rejects duplicates and out-of-range phases,
checks pulse length and finite/bounded taps, and limits JSON size to 128 KiB.
Those structural checks do not prove that an arbitrary supplied profile matches
a real modem. Reproducible tests use SHA-256-pinned upstream measurements,
extracted without executing downloaded source. Cached measurement sources
retain their own attribution/license metadata and are not embedded in Rust.

The phase index $k_{b,t}$ selects one of 128 equally spaced directions
around a circle. For example, index 32 represents a quarter turn. The
transmitter and receiver use opposite rotations so they can compare
successive DATA symbols after accounting for the known phase pattern.

\landscapediagram{\begin{tikzpicture}[x=1.02cm,y=1cm,font=\small,
  axis/.style={ovgray,-{Stealth[length=2mm]}},
  data/.style={ovblue,fill=ovblue},
  pilot/.style={ovgold!60!black,fill=ovgold},
  reference/.style={ovteal!70!black,fill=ovteal}]
  \node[anchor=west,font=\bfseries] at (-2.7,6.15) {BW500 level-4 symbol centers};
  \node[anchor=west] at (-2.7,5.55) {$n_{b,j}=4608+512j+256b$, with band index $b\in\{0,1\}$ and column $0\le j\le392$};
  \foreach \x in {0,1,...,17}
    \draw[ovgray!20] (\x,1.05) -- (\x,4.95);

  \draw[axis] (0,3.9) -- (17.65,3.9);
  \draw[axis] (0,2.45) -- (17.65,2.45);
  \node[align=right,anchor=east] at (-0.35,3.9) {Lower band\\1,350 Hz};
  \node[align=right,anchor=east] at (-0.35,2.45) {Upper band\\1,650 Hz};
  \draw[reference] (-0.09,3.81) rectangle ++(0.18,0.18);
  \draw[reference] (0.91,2.36) rectangle ++(0.18,0.18);
  \foreach \j in {1,...,8} {
    \fill[data] ({2*\j},3.9) circle[radius=2.4pt];
    \ifnum\j=7
      \draw[pilot] (15,2.59) -- (15.14,2.45) -- (15,2.31) -- (14.86,2.45) -- cycle;
    \else
      \fill[data] ({2*\j+1},2.45) circle[radius=2.4pt];
    \fi
  }
  \node[above,font=\scriptsize] at (0,4.11) {$j=0$};
  \node[above,font=\scriptsize] at (2,4.11) {$j=1$};
  \node[above,font=\scriptsize] at (15,2.59) {$j=7$, first upper pilot};
  \draw[ovblue,{Stealth[length=1.4mm]}-{Stealth[length=1.4mm]}]
    (0,4.7) -- (2,4.7) node[midway,above,font=\scriptsize] {512-sample hop};
  \draw[ovteal!70!black,{Stealth[length=1.4mm]}-{Stealth[length=1.4mm]}]
    (0,3.15) -- (1,3.15) node[midway,above,font=\scriptsize] {256};

  \draw[axis] (0,1.05) -- (17.65,1.05);
  \foreach \x/\sample in {0/0,2/512,4/1024,6/1536,8/2048,10/2560,12/3072,14/3584,16/4096}
    \draw[ovgray] (\x,1.05) -- ++(0,-0.11) node[below,font=\scriptsize] {\sample};
  \node[anchor=west] at (0,0.34) {Samples relative to the lower-band reference, $n-4608$};

  \draw[reference] (0,-0.55) rectangle ++(0.18,0.18);
  \node[anchor=west] at (0.32,-0.46) {Reference};
  \fill[data] (3.05,-0.46) circle[radius=2.4pt];
  \node[anchor=west] at (3.27,-0.46) {Data};
  \draw[pilot] (5.1,-0.32) -- (5.24,-0.46) -- (5.1,-0.60) -- (4.96,-0.46) -- cycle;
  \node[anchor=west] at (5.34,-0.46) {Pilot};
  \node[anchor=west,font=\scriptsize] at (7.0,-0.46)
    {18 pilots per band. The first lower-band pilot is at $j=24$.};
\end{tikzpicture}
}{BW500 level-4 symbol centers on the 48 kHz sample axis. The plotted excerpt uses the implemented pilot allocation. Each marker locates the center of a 513-tap pulse extending 256 samples on either side.}{fig:dsp-level4-symbol-timing}

\subsection{Level-4 transmit synthesis}
After whitening, turbo encoding, and channel gather, coded bits fill band 0's
data positions in time order, then band 1's. With initial polarity
$v_{b,0}=+1$, a coded 1 negates polarity, a coded 0 preserves it, and a pilot
repeats the preceding polarity. Each complex emitted cell is
$c_{b,t}=v_{b,t}\overline{G_{b,t}}$. The real audio is a sum of pulse-shaped
carriers:
\[
 x[n]=\frac14\sum_{b=0}^{1}\sum_t p[n-n_{b,t}]
       \Re\{c_{b,t}\exp(j2\pi f_b n/48000)\},
\]
where $n_{b,t}=9H+tH+bH/2$, $f_b\in\{1350,1650\}$, and $p$ has 513 taps.
To compute one output sample $x[n]$, center each symbol's pulse on
$n_{b,t}$, evaluate the carrier at sample $n$, and add the real
contributions from both bands. The pulse determines how much each nearby
symbol contributes. The center formula allows nine hops before the
first reference, advances one hop per column, and offsets the upper
band by half a hop.

The factor $1/4$ is the current bench amplitude scale, not a calibrated radio
power command. A waveform's audio amplitude does not replace CAT power limits,
route qualification, or PTT authority.

The synthetic buffer contains $9H+393H+1=205825$ samples, approximately
4.288 s at 48 kHz. The initial slots do not contain a generated randomized
acquisition lead. This output length is a local implementation fact and must
not be substituted for a measured reference modem's PTT interval, framing
duration, or turnaround requirement. Actual device playback and RF keying are
not performed by the file-only synthesis tool.

\subsection{Level-4 timing lock, demodulation, and validation}
\modemcite{proakis,meyr}The receiver mixes the real input down at each nominal carrier using
$x[n]\exp(-j2\pi f_b n/48000)$. At a trial symbol center it multiplies the
513 surrounding complex samples by the measured real pulse and sums them.
It then multiplies by $G_{b,t}$ to remove the known phase reference. If the
result is $z_{b,t}$, the differential statistic is
\[
 d_{b,t}=\Re\{z_{b,t}\overline{z_{b,t-1}}\}.
\]
This dot product compares consecutive symbols after removing the known
reference phase. It tends to be positive when their polarities agree
and negative when they are opposite. A DATA bit 1 flips polarity,
so the receiver negates this statistic to give a positive score for 1.

A correct pilot transition should preserve polarity. Timing search maximizes
the average normalized differential dot product over 36 pilot positions,
first in 32-sample steps and then in two-sample steps around the best epoch.
A score below 0.65 is rejected. The search is bounded by the candidate, frame
geometry, available samples, and one-hop end tolerance; it does not scan an
unbounded recording for a matching payload.

Data positions produce $-d_{b,t}$ so positive metrics favor a transmitted 1.
The 748 values are scaled by six divided by their mean absolute magnitude,
then passed through channel scatter and the shared turbo/de-whitening/CRC
chain. This normalization is not an estimate of a physical noise variance.
No valid CRC means no frame. A valid frame is interpreted in protocol context:
a setup frame can identify a caller; a DATA container can carry payload and
control; neither is an unconditional link-up event.

\landscapeplantuml{build/diagrams/Linkwise_DSP_Level4}{Implemented level-4 transmit and receive transforms. Measured reference phases and pulse taps are external profile inputs; the single-frame synthetic round trip does not qualify reference interoperability.}{fig:dsp-level4}

The level-4 decoder currently allocates complex mixed buffers and performs
repeated direct matched-pulse sums. It is a bench implementation, not a
profiled real-time DSP scheduler. An optimized runtime may cache oscillators,
reference terms, pulse weights, and mode tables, and bound search work per
candidate, but must preserve sample epochs, metric sign, permutation order,
and CRC outcomes against golden fixtures before replacing these kernels.


\subsection{Proposed production DSP extensions and work budgets}
The following is a proposed engineering path, not behavior implemented by the
current decoder. A production receiver needs explicit acquisition/tracking
states, measured latency budgets, and loss-of-lock rules rather than extending
the bench thresholds by assumption. Mode tables, pulse taps, and reference
terms should be validated and cached before a session; source loading,
allocation-heavy analysis, persistence, and model calls stay outside the
capture and playback deadlines.

\begin{enumerate}[leftmargin=1.6em]
  \item \textbf{Front-end conditioning and health.} Measure route noise,
        clipping, DC offset, stream discontinuities, and usable passband.
        Select channel filters and gain normalization from qualified profiles.
        An adaptive candidate detector needs a recorded noise estimator,
        hysteresis, duration limits, and a false-alarm/miss test population.
        Gain changes must not erase sample epochs or turn audio energy into
        protocol evidence.
  \item \textbf{Coarse acquisition.} Search a bounded frequency/timing bank
        against references before attempting expensive FEC. Record competing
        mode scores and ambiguity rather than treating the first energetic
        interval as a known modem. Reuse selected-bin correlations or an FFT
        bank only after latency measurements establish the benefit and fixture
        tests preserve the selected bins and normalized scores.
  \item \textbf{Carrier tracking.} After correct reference de-rotation and
        under a locally constant channel, repeated pilot transitions provide
        a candidate residual-offset estimator:
        \[
          \widehat{\Delta f}=\frac{48000}{2\pi H}
          \arg\!\left(\sum_{(b,t)\in\mathcal P}
          w_{b,t}z_{b,t}\overline{z_{b,t-1}}\right).
        \]
        Each pilot pair measures the phase change over one symbol hop.
        The set $\mathcal P$ identifies the pilot positions and the
        weights $w_{b,t}$ determine how much each pair contributes.
        The angle of the combined result gives radians per hop.
        Multiplying by $48000/(2\pi H)$ converts that angle to hertz.
        For example, a quarter-turn of phase per 512-sample hop
        corresponds to about 23.44 Hz.

        This is a proposed pilot estimator, not a performance result. At
        $H=512$ its phase observation repeats every 93.75 Hz, requiring coarse
        branch selection and outlier/low-energy rejection. Fading and channel
        variation must be included in qualification before closing a loop.
  \item \textbf{Sample-clock and channel tracking.} Track fractional symbol
        epochs and their drift over long bursts; qualify interpolation or
        resampling with known offsets rather than slipping whole samples
        silently. Assess pulse distortion and multipath-induced intersymbol
        or interband interference. Any adaptive equalizer needs training,
        convergence, divergence, and reset criteria against recorded HF
        channel cases. A capture gap invalidates continuity and tracking state.
  \item \textbf{Metric and deadline qualification.} Compare heuristic
        power/differential metrics with noise-aware likelihood models; validate
        sign, clipping, erasures, and decoder response before changing scales.
        Bound candidate search, queue occupancy, FEC iteration count, and
        cancellation latency. Measure worst-case and high-percentile service
        times on each target under simultaneous capture, decode, client I/O,
        and audit load. Drop optional diagnostics before compromising sample
        continuity; if continuity or a safety deadline cannot be maintained,
        terminate the affected stream/session explicitly.
\end{enumerate}

The current turbo decoder performs forward/backward passes on two eight-state
trellises, so work scales linearly with information size and iteration count;
forward metric storage is also linear in block size. This algorithmic bound
does not establish a real-time scheduling bound. Index acquisition repeatedly
correlates reference windows, and level-4 trial epochs repeatedly evaluate
matched pulses; both can dominate a failed candidate. Kernel benchmarks must
therefore include failed acquisition and failed CRC/FEC paths, not only the
fast hard-decision success case.

A DSP qualification record should identify source audio/hash, route and sample
format, mode/profile hash, impairment model and seed, expected frame/payload,
acquisition epoch, frequency/clock estimates, decoder iterations, terminal
CRC result, queue high-water marks, and elapsed work. BER requires known
reference bits; a CRC failure alone does not measure BER. Fixture-defined
frame-error rates, energy-candidate rejection rates, and link/session outcomes
need separate denominators. Actual output presentation latency and audio/PTT
lead/tail behavior require device-loopback measurements in addition to these
software records.

\subsection{Control bursts, caller setup, and handshake evidence}
Control bursts use a separate 2,048-sample tone grid at 48 kHz. Request and
response tones are deterministic functions of the known called station;
their payload uses the BW500 even-bin alphabet 50--76 while fixed preambles
may use other bins. The request contains ten preamble plus 31 payload symbols;
the response contains eight plus 15; the level-4 confirm contains one plus 15.
Connected ACKs contain eleven two-tone symbols, with a fixed four-symbol
preamble and a caller-keyed tail. The control renderer uses a 32-sample overlap
and cosine edge shaping. These control bursts do not pass through the DATA
turbo/interleaver chain.

Recognizers compare the complete expected keyed sequence near the estimated
burst start. Single-tone matching requires an average normalized tone score
of at least 0.8 and a weakest-symbol score of at least 0.30; connected-ACK
matching requires total pair score of at least 0.8 per symbol and at least
0.15 for each individual expected tone. These are bench thresholds, not
field-qualified receiver sensitivities. The complete ACK tail is checked;
shared preamble energy alone cannot establish acceptance. Request/response
matching cannot discover an unknown caller because those bursts are keyed to
the called station. Caller identity arrives in the CRC-protected setup frame.

The setup codec packs up to six base callsign characters into six-bit symbols,
retains an SSID in byte 5, records capability/request fields, includes the high
byte of the full-callsign CRC, and finishes with a frame CRC. Current native
audio establishment is restricted to level 4. Its state machine expects request,
response, setup, connected ACK, then confirm; wrong-phase or wrong-peer bursts
cannot advance it. Unique local transmit tokens reject stale completion
callbacks. Establishment is emitted only after the relevant confirm is
received or actually transmitted, rather than when a waveform is merely queued.

The current two-role integration test supplies transmit completion after the
synthetic audio has been exchanged. A production device owner must provide
that event only after actual presentation and PTT release. Software buffer
consumption is not proof that the final sample reached the output device.
Clocked playback, transmit/receive turnaround, and qualified PTT readback
remain required integration work.

\subsection{DATA framing, ARQ ownership, and ambiguous answers}
The DATA container codec is separate from waveform decoding. At level 4 the
body has 44 bytes, at level 3 it has 35, and the final body byte is a control
field. A full frame uses the preceding bytes for payload. The supported short
form places payload, a \texttt{0x14} marker, the caller CRC high byte, zero
fill, a next-receive-level byte, and \texttt{0x82} before the CRC. Unsupported
or ambiguous forms are rejected. The implemented short builder reserves four
body bytes for this trailer; additional published short/compressed forms are
not silently treated as supported.

The BW500 decoder also accepts the measured level-4 \texttt{0x02}
control form. In setup, it validates the CRC, caller identity and
requested level. In DATA, it requires the caller-keyed trailer and valid
padding and rate fields. Stock Winlink Express used this form to send
the B2F end-of-proposals bytes after storing the received message.

Tests reconstruct the documented 200-byte delivery from four level-4 payloads
of 43 bytes and one closing level-3 payload of 28 bytes, using the documented
full fields \texttt{95}, \texttt{91}, \texttt{8d}, and \texttt{81}. This verifies
the container interpretation for that example, not a general sender sequencing
law. The full-frame field remains a supplied parameter because its evolving
wire semantics need independent qualification. Appendix~\ref{app:bw500-wire}
works through these bytes, their checksums, and the session exchange.
Compression TEXT/FILES,
delivery-level repacketization, and negotiated rate changes are not completed
by the container codec.

The event-level ARQ sender owns bounded queued blocks and exactly one
outstanding block. DATA transmit completion opens a receive deadline. A
validated delivered reply retires that block once; a validated not-delivered
reply permits a retry within budget. A missing reply is ambiguous: the peer
may have accepted DATA while its answer was lost. The timeout therefore emits
a state-query action rather than blindly repeating DATA. Query transmission
has its own completion token and watchdog. Query/retry exhaustion, transmit
failure, or missing/late playback completion faults the sender and preserves
unresolved bytes. Explicit abort discards the queue.

These local tokens bind callbacks to an outstanding transmission/window; they
are not transmitted sequence numbers or a solution to on-air duplicate
suppression. The station/native session owner must connect qualified reply
recognition, query encoding, turn ownership, negotiated rate, and byte
accounting. The sender's queue currently holds supplied blocks; its byte count
must not be exposed as Winlink delivery progress or raw-host \texttt{BUFFER}
semantics without that integration.

The virtual host issues keyed state queries after missing BW500
intermediate and closing acknowledgments. A positive answer retires the
acknowledged bytes. A negative answer can trigger a level-3 retry of the
outstanding DATA, preserving any acknowledged prefix. With crossed turn
requests, Linkwise can grant a repeated stock caller request and resume
queued DATA after the caller hands back the turn. A short continuation
can retain level 3 after a level-3 close. A host \texttt{DISCONNECT}
received while DATA is queued waits for the queue to drain before
sending the disconnect waveform.

For the wide modes, the virtual host tracks the outstanding waveform and
the byte offset at the active rate. It can retry the unacknowledged suffix
at level 3 after a keyed rejection, use level 2 after a missing-DATA query
response, and return to level 4 when the peer requests the supported
increase. These transitions cover measured frame layouts. General rate
negotiation and repeatable recovery across the qualification matrix
remain open.

\landscapeplantuml{build/diagrams/Linkwise_Modem_Recovery}{Implemented event-level ARQ sender behavior and remaining wire integration. Missing answers request state; only a validated negative reply permits DATA retransmission.}{fig:modem-recovery}

\subsection{Execution isolation, keying, and failure containment}
The proposed runtime separates four execution responsibilities. The input
callback only converts and queues audio; DSP workers detect, synchronize,
decode, and synthesize outside that callback; one native session actor owns
protocol transitions, outstanding bytes, and monotonic deadlines; the station
execution service owns the radio lease and client/modem operation. The PTT
arbiter and watchdog retain independent release authority. MCP queries read
snapshots and append commands to bounded application queues; slow model
responses, JSON serialization, audit storage, and provider fetches never run
inside an audio callback or protocol deadline.

A burst request must identify the session, active lease, configuration
revision, permitted route, and bounded duration. The arbiter validates those
against policy and interlocks before key-up, arms release limits, and uses a
qualified driver. Playback completion and release observation are separate
facts. Timeout, device loss, capture discontinuity, lease loss, abort, or an
unconfirmable release cannot be converted to a successful session. A release
request still proceeds when audit persistence fails; the fault is latched and
reported separately. The physical emergency stop remains independent of
software execution and of either AI host.

\landscapeplantuml{build/diagrams/Linkwise_Modem_Keying}{Proposed local playback, keying, completion, and watchdog sequence below an asynchronous agent request. This integration is not yet implemented; no MCP round trip is allowed inside a symbol or ACK deadline.}{fig:modem-keying}

\subsection{Verification matrix and implementation exit gates}
The 9 November run of \path{tools/test-modem.py --audio-device} passed 41 tests:
37 library tests, one replay integration test covering five independently
published level-1/2 excerpts, and three synthetic level-4/audio integration
tests. The reference clips include a two-frame burst and retain upstream
provenance and SHA-256 checks. The level-4 tests cover setup and DATA round
trips, local two-role establishment, chunked streaming, and corruption
rejection. Formatting and Clippy checks also passed. These are dated evidence,
not a promise that subsequent source changes have identical results.

On 10 November the fixture-enabled suite passed 68 tests on macOS
with the audio feature enabled. The pilot acquisition regression
covered consecutive frames in continuous interference, chunk independence,
and duplicate suppression. Formatting and Clippy passed. The offline
receive checks used the full demodulated recordings and supplied no burst
boundaries. Capture settings, hashes, and reproduction commands are in the
10 November entry of \path{lab-notes/lab-notes.tex}.

The protocol notes and reference measurements are pinned to upstream revision
\path{952f79245240132f4af209afbad1f4b2ba256613}; implementation details should
be checked against the named Rust modules and \path{rust/docs/modem.md}.
The notes, fixtures, and measured sources have distinct attribution/license
metadata; the cached measurement inputs are not proprietary vendor assets
embedded in this repository. The public protocol/reference source is listed
in Reference~\ref{ref:modem-protocol}. Source-level traceability does not
replace empirical qualification.

\paragraph{Deterministic waveform regression and application delivery.}
The 12 November 2026 BW500 suite passes packed PRBS9, PRBS15, alternating
bits, and all-zero and all-one payloads through Linkwise's real-valued
48~kHz waveform path. Levels 1--3 recover the exact DATA frames through
attenuation and a fixed 1~ms echo. A separate level-4 test reassembles
the 64-byte PRBS9 fixture from full and closing DATA frames using the
measured waveform profile. The virtual daemon bench also recovers
PRBS9 payloads at BW500, BW2300, and BW2750, with controlled DATA
corruption or acknowledgment loss in selected cases.

On the mixed BW500 bench, both Winlink Express peers sent and received
an email with a PRBS9 attachment in each of two sessions. Linkwise and
stock VARA HF 4.9.1 each initiated one session. Both peers disconnected
cleanly and the four received attachments matched their sources byte for
byte and by SHA-256. Separate stock-audio replays cover the initial caller
poll and rate-down recovery in both caller roles. These measurements
cover frame decoding and delivery after ARQ. Measuring raw bit errors
requires aligned receiver decisions, and assessing fading tolerance
requires a calibrated channel model. Fixture construction, commands, and
run scope appear in the 12 November entry of \path{lab-notes/lab-notes.tex}.

\begin{longtable}{@{}p{0.24\textwidth}p{0.70\textwidth}@{}}
\caption{Modem evidence and remaining qualification gates.}\label{tab:modem-evidence-gates}\\
\toprule
\textbf{Gate} & \textbf{Evidence and remaining checks}\\
\midrule
\endhead
Codec correctness & Independent vectors for whitening, permutations, FEC tails,
metric sign, frame fields and mode geometry; expand beyond self-generated
round trips and the level-1/2 corpus.\\
Off-air receive & One CRC-valid BW500 level-2 frame from the full Lake
Merwin KiwiSDR recording on 7.105 MHz USB. The published off-air exchange
yielded three CRC-valid frames (levels 2, 3, 2). Both replays reproduced
previously verified wire-body hashes.\\
Acquisition controls & Two stock-derived frames in recorded Kiwi noise
decoded automatically at +6 dB. The Maui noise-only replay emitted zero
frames. These checks cover the tested recordings.\\
Receiver robustness & Measure live frame recovery on known transmissions.
Characterize offset, fading, interference, clipping, sample-clock error,
false alarms, and processing latency.\\
Reference acceptance & With stock VARA HF 4.9.1, the 11 November frozen
build \texttt{6a4dc89e} passed one BW2300 and two BW2750 sender-recovery
runs with 412 exact bytes, two-ended disconnect and a same-socket
reconnect. Other runs stalled during setup or retry. The three-pass gate
remains open.\\
Session reliability & The frozen build passed two final-ACK-loss runs and
one intermediate-ACK-loss run at BW500, with 406 exact bytes and two-ended
disconnect in each pass. Setup and DATA timeouts interrupted the sequences.
The earlier three-pass results apply to their recorded builds.\\
Station safety & Qualify actual device presentation, turnaround, exclusive
lease, permit expiry, duty limits, watchdog release and independent stop on
an inhibited or conducted setup before operational RF use.\\
Host and Winlink & Mixed BW500 Pat runs passed in both directions with a
matching attachment SHA-256, exit status zero and two-ended disconnect.
Mixed Winlink Express 1.8.6.0 and Pat 1.0.1 peer-to-peer transfers stored
attachments with matching SHA-256 in both recipient mailboxes. The
Pat-to-Express sender-side close and gateway/CMS qualification remain open.
Two Winlink Express BW500 sessions with Linkwise and stock VARA verified
two-way PRBS9 attachment delivery and clean disconnect with each modem
as caller.\\
\bottomrule
\end{longtable}

The earlier i386 Debian/Wine bench on the ARM64 Mac stopped after an
illegal-instruction exception during COM registration and a WineDbg
attach error. The later Windows UTM bench provided a working stock peer.
The November 11 lab chronicle records the stock Windows UTM runs, the
Express message export, and installer checks. Its evidence manifest links
the measured results to the captured audio and command transcripts.
The repeat counts, build hashes, and failure captures are indexed in
\path{docs/testing/vara-interoperability.md}.
Physical transmit control and operational RF qualification remain open.
